% !TeX root = ../../Book1.tex
\section{可整流流、积分流与几何运算}
\label{b1:sec:D103}

正规性控制质量和边界质量，却没有规定质量应集中在什么集合上。
若要描述带重数的曲面，应加入可整流性；
若希望一张曲面在极限中仍由整数层组成，还要加入整数重数。
这几个要求逐层增加，不能把“有面积”“有边界”和“有整数层数”当成同一条件。

\subsection{承载集、重数和取向}
\label{b1:subsec:D10301}

\begin{definition}\label{b1:def:app-d-rectifiable-current}
设 \(1\leq k\leq n\)。若一个流可以写成
\[
 T(\omega)=\int_E\theta(x)
                 \langle\omega(x),\tau(x)\rangle\,d\mathcal H^k(x),
\]
其中 \(E\) 可数 \(k\) 维可整流，\(\theta>0\) 可测，
\(\theta\mathcal H^k\llcorner E\) 局部有限，
且 \(\tau(x)\) 几乎处处是近似切空间的可测单位可分解 \(k\) 向量，
则称 \(T\) 为%
\term[kezhengliu-liu]{可整流流}[rectifiable current]，记作 \(T=[E,\theta,\tau]\)。
若另有 \(\theta(x)\in\mathbb N\) 几乎处处成立，则称其具有整数重数。
\end{definition}

本书在“可整流流”中容许实重数；需整数时明确写出。
Simon 的《Lectures on Geometric Measure Theory》%
\cite[\S27]{Simon1983Lectures}专门讨论 integer multiplicity rectifiable currents，
并说明其与 Federer--Fleming 用语的对应。
查阅其他资料时，不能仅凭 rectifiable 一词判断是否已经包含整数条件。

\ProseParagraphBreak

方向的符号由 \(\tau\) 携带，故取 \(\theta>0\) 不损失一般性。
也可以固定一套可测取向，允许 \(\theta\) 有正负号；两种表示给出相同的流。
由\cref{b1:prop:app-d-mass-representation}以及可分解向量的质量范数，
\[
 \|T\|=\theta\mathcal H^k\llcorner E,\quad
 \mathbf M(T)=\int_E\theta\,d\mathcal H^k.
\]
这里的重数并不是把集合真的复制到别处，而是在同一位置把积分计入多次。

\ProseParagraphBreak

零维的约定单独说明：局部整数重数零维流是
\(\sum_iq_i\delta_{x_i}\)，其中 \(q_i\in\mathbb Z\setminus\{0\}\)，
每个紧集内 \(\sum_{x_i\in K}|q_i|<\infty\)。
因每个非零整数的绝对值至少为一，这等价于每个紧集只含有限个非零点。
全域有限质量时，非零点总数有限。

\begin{definition}\label{b1:def:app-d-integral-current}
若 \(T\) 及 \(\partial T\) 都具有整数重数可整流表示，且
\[
 \mathbf M(T)+\mathbf M(\partial T)<\infty,
\]
则称 \(T\) 为%
\term[jifen-liu]{积分流}[integral current]。
我们把 \(\mathbb R^n\) 上全体 \(k\) 维积分流构成的集合记为 \(\mathbf I_k(\mathbb R^n)\)。
\symidx{\(\mathbf I_k(\mathbb R^n)\)：有限质量积分流}
只在内部紧集上要求相应有限性时，使用“局部积分流”。
\end{definition}

这一定义已经把边界可整流性列入条件，不依赖尚未证明的边界可整流定理。
在更深入的理论中，整数重数可整流流只要边界局部质量有限，
边界便自动整数重数可整流；Simon 的 30.3 THEOREM\cite[30.3 THEOREM]{Simon1983Lectures}给出这一结果。
本附录的紧性和变分问题使用上述明示边界条件的定义。

\begin{proposition}\label{b1:prop:app-d-integral-algebra}
对每个 \(0\leq k\leq n\)，\(\mathbf I_k(\mathbb R^n)\)
对加法、取负和整数倍封闭。
若 \(1\leq k\leq n\) 且 \(T\in\mathbf I_k(\mathbb R^n)\)，则
\(\partial T\in\mathbf I_{k-1}(\mathbb R^n)\)。
\end{proposition}
\begin{proof}
零维情形由有限整数系数点质量的表示直接得到。
以下设 \(k\geq1\)。在两个可整流承载集的交集上，
它们的近似切空间对 \(\mathcal H^k\)-几乎处处的点相同。
这可由可整流集的光滑片分解得到：两张同维光滑片若切空间不同，
其交集在该点附近的 \(k\) 维面积密度为零。
因此在公共部分，两套单位定向相同或相反，重数相加或相减后仍是整数。
删去重数为零的部分，并把负号吸收到定向中，便得到和的整数可整流表示。
三角不等式给出质量有限性，边界之和也有同样性质。
取负、整数倍相同处理。
最后 \(\partial^2T=0\)，所以 \(\partial T\) 的边界是零流，
由积分流的定义可知 \(\partial T\in\mathbf I_{k-1}(\mathbb R^n)\)。
\end{proof}

有限个有向 \(k\) 维单纯形的整数线性组合给出最直接的例子。
称这样的流为整数系数多面体流；允许把单纯形换成有向多面体面，经过有限剖分所得类别不变。
相邻面的公共边界按取向相消，所以一个剖分好的平面区域不会在每条内部剖分线上产生额外边界。
在这种对象上，流的加法就是有向链的加法。

\subsection{有限周长集合给出的积分流}
\label{b1:subsec:D10302}

令 \(E\subset\mathbb R^n\) 可测，定义最高维流
\[
 [E](f\,dx^1\wedge\cdots\wedge dx^n)=\int_E f\,dx.
\]
它的质量等于 \(|E|\)。
最高维切空间就是整个 \(\mathbb R^n\)，所以这个表示自然具有重数一和标准取向。

\begin{proposition}\label{b1:prop:app-d-perimeter-current}
若 \(E\) 具有有限体积和有限周长，则 \([E]\) 为积分流，且
\[
 \mathbf M([E])=|E|,\quad
 \|\partial[E]\|=\mathcal H^{n-1}\llcorner\partial^*E.
\]
边界取向 \(\tau_E\) 由
\[
 \nu_E\wedge\tau_E=e_1\wedge\cdots\wedge e_n
\]
确定，其中 \(\nu_E\) 为集合的测度论外法向。
\end{proposition}
\begin{proof}
每个 \((n-1)\) 次形式都唯一写成
\[
 \eta=\sum_{i=1}^n(-1)^{i-1}X_i\,
        dx^1\wedge\cdots\wedge\widehat{dx^i}\wedge\cdots\wedge dx^n
\]
其中 \(X\) 为光滑向量场，帽号表示删去该项。
直接外微分得到 \(d\eta=(\operatorname{div}X)\,dx^1\wedge\cdots\wedge dx^n\)。
由\cref{b1:thm:gauss-green-finite-perimeter}，
\[
 \partial[E](\eta)
   =\int_E\operatorname{div}X\,dx
   =\int_{\partial^*E}X\cdot\nu_E\,d\mathcal H^{n-1}.
\]
上述取向满足
\(\langle\eta,\tau_E\rangle=X\cdot\nu_E\)，故
\(\partial[E]=[\partial^*E,1,\tau_E]\)。
由 De Giorgi 结构定理，\(\partial^*E\) 可数 \((n-1)\) 维可整流，
其面积等于周长；两个质量均有限。
\end{proof}

从这个例子可看出，有界变差函数理论已经包含一种重要的流。
集合的拓扑边界可能因修改零测集而完全改变，流 \([E]\) 及其边界却不会改变。
边界也不应写成整个拓扑边界上的面积积分，正确承载集是约化边界。

\subsection{光滑映射下的推前}
\label{b1:subsec:D10303}

设 \(f:U\to V\) 光滑。
外代数的拉回给出
\[
 (f^*\omega)_x(\xi)
   =\omega_{f(x)}\bigl((\bigwedge\nolimits^kDf_x)\xi\bigr).
\]
若每个目标紧集的逆像与 \(\operatorname{supp}T\) 相交后仍紧，则定义
\[
 (f_\#T)(\omega)=T(f^*\omega).
\]
当拉回不在全域紧支集时，右端在它与 \(\operatorname{supp}T\) 相交的紧集附近乘一个等于一的截断。
不同截断的差在支撑附近为零，因而给出同一数值。
本附录后面主要使用紧支集流，此时连续映射自动满足所需的支撑条件。

\begin{proposition}\label{b1:prop:app-d-pushforward}
在上述支撑条件下，推前是流，且
\[
 \partial(f_\#T)=f_\#(\partial T).
\]
若 \(T\) 质量有限且 \(Df\) 在其支撑上有界，则
\[
 \mathbf M(f_\#T)\leq
   \left(\sup_{\operatorname{supp}T}\|Df\|_{\mathrm{op}}\right)^k
          \mathbf M(T).
\]
若 \(T\) 为整数重数可整流流，则 \(f_\#T\) 也具有整数重数可整流表示；
若 \(T\) 为紧支集积分流，则 \(f_\#T\) 为积分流。
\end{proposition}
\begin{proof}
链式法则保证拉回在固定紧支集上的各阶导数受控，所以推前连续。
逐坐标计算给出 \(d(f^*\eta)=f^*(d\eta)\)：
一阶系数满足通常链式法则，二阶偏导项因反对称性相消。
于是
\[
 \partial(f_\#T)(\eta)=T(f^*d\eta)
       =T(d f^*\eta)=f_\#(\partial T)(\eta).
\]
对单位可分解向量 \(\xi\)，其像仍可分解，且
\(|(\bigwedge^kDf)\xi|\leq\|Df\|_{\mathrm{op}}^k\)。
由此
\(\|f^*\omega\|_{\mathrm{co}}\leq
\|Df\|_{\mathrm{op}}^k\cdot\|\omega\circ f\|_{\mathrm{co}}\)，给出质量估计。

若 \(T=[E,\theta,\tau]\)，则
\[
 f_\#T(\omega)
   =\int_E\theta(x)
       \left\langle\omega(f(x)),(\bigwedge\nolimits^kd^Ef_x)\tau(x)\right\rangle
          \,d\mathcal H^k(x).
\]
对切向 Jacobi 因子正的部分应用可整流集上的面积公式。
在几乎每个像点，所有非退化原像切空间的像等于该像集的近似切空间；
各项取向因此只差一个正负号。
选定像集的一套可测取向后，其重数为各原像整数重数的带符号和。
质量估计保证这些绝对值的和几乎处处有限，所以该和仍为整数。
切向 Jacobi 因子为零的部分不给积分贡献。
对 \(T\) 及 \(\partial T\) 分别使用该结论，再结合边界交换式，即得最后一句。
\end{proof}

像的几何面积不能总用带符号重数的绝对值之前的求和来代替。
两片反向曲面映到同一位置时会相消，故质量估计一般只是“不超过”。
Lipschitz 映射也有推前理论，但一般流的测试形式不能直接拉回成光滑形式。
本节先使用光滑版本；在建立积分流紧性后，%
\cref{b1:prop:app-d-lipschitz-pushforward}将利用同伦控制完成
紧支集积分流上的 Lipschitz 延拓。

\subsection{乘积与同伦公式}
\label{b1:subsec:D10304}

设 \(T\) 为紧支集正规 \(k\) 维流。
在 \(\mathbb R\times\mathbb R^n\) 中，把测试 \((k+1)\) 次形式写成
\(\Omega=dt\wedge\alpha(t,x)+\beta(t,x)\)，
其中 \(\alpha\) 不含 \(dt\)，\(\beta\) 也不含 \(dt\)。
定义
\[
 ([0,1]\times T)(\Omega)=\int_0^1T\bigl(\alpha(t,\cdot)\bigr)\,dt.
\]
该积分的绝对值受质量界控制，且支撑包含在
\([0,1]\times\operatorname{supp}T\) 内，因此定义了有限质量流。

\begin{proposition}\label{b1:prop:app-d-homotopy}
设 \(h:[0,1]\times\mathbb R^n\to\mathbb R^N\)
在该闭条带的邻域上光滑，并记 \(h_t(x)=h(t,x)\)。
对紧支集正规流 \(T\)，有
\[
 \partial h_\#([0,1]\times T)
  =(h_1)_\#T-(h_0)_\#T-h_\#([0,1]\times\partial T).
\]
\end{proposition}
\begin{proof}
先验证乘积边界。
把一个 \(k\) 次测试形式写成
\(\eta=dt\wedge\gamma(t,x)+\alpha(t,x)\)。
其外微分中含 \(dt\) 的部分为
\[
 dt\wedge(\partial_t\alpha-d_x\gamma).
\]
所以
\[
 \begin{split}
 \partial([0,1]\times T)(\eta)
 &=\int_0^1 T(\partial_t\alpha)\,dt
     -\int_0^1\partial T(\gamma)\,dt\\
 &=T(\alpha(1,\cdot))-T(\alpha(0,\cdot))
     -([0,1]\times\partial T)(\eta).
 \end{split}
\]
第一项积分使用一维微积分基本定理；通过质量表示可直接把 \(T\) 与积分交换。
然后应用推前与边界交换的结论，即得公式。
\end{proof}

这条公式说明，两个推前之间的差不一定本身就是一条边界：
当原流有边界时，边界移动扫出的侧面还要计入。
在估计平坦范数时遗漏这项，就会错误地只用 \(\mathbf M(T)\) 控制所有平移。

\ProseParagraphBreak

特别地，令 \(h(t,x)=a+t(x-a)\)。
若 \(T\) 是紧支集积分 \(k\) 维流且 \(k\geq1\)，则
\((h_0)_\#T=0\)；
若又有 \(\partial T=0\)，便得到一个填充
\[
 C_aT=h_\#([0,1]\times T),\quad \partial C_aT=T.
\]
乘积具有整数重数可整流表示，光滑推前保持该性质，因此 \(C_aT\) 为积分流。
若 \(\operatorname{supp}T\subset\overline{B(a,R)}\)，
其切向外积在时间 \(t\) 的长度至多为 \(R t^k\)，从而
\[
 \mathbf M(C_aT)\leq\frac{R}{k+1}\,\mathbf M(T).
\]
在零维时，常值推前等于总系数乘 \(\delta_a\)，所以只有总系数为零的零维流才能这样填充。
例如 \(\delta_b-\delta_c\) 有线段填充，而单个 \(\delta_b\) 没有紧支集的一维填充。
